\SourceLeaf{355}

\begin{center}
\large\bfseries GEOMETRIC ASPECT OF TERNARY FORMS.
\end{center}



Let a point $(0)$ in space be taken as the origin. The displacements from there to three
other points $\boldsymbol{P}$, $\boldsymbol{P}'$, $\boldsymbol{P}''$, which do not lie
in a plane with it, are respectively $t$, $t'$, $t''$. Whenever no confusion can arise,
the points $\boldsymbol{P}$, $\boldsymbol{P}'$, $\boldsymbol{P}''$ themselves may be denoted by $(t)$,
$(t')$, $(t'')$.

Further, let $(t,t')$ in general be the product of the lengths (shortest distances) of the two
lines $t$, $t'$ and of the cosine of their inclination, etc.

In general, we obtain $(\alpha t + \alpha' t' + \alpha'' t'' + .\,.\,.,\ \beta u + \beta' u' +
\beta'' u'' + .\,.\,.)$ by carrying out the multiplication
\[
(\alpha t + \alpha' t' + \alpha'' t'' + .\,.\,.) \times (\beta u + \beta' u' + \beta'' u'' + .\,.\,.)
\]
and replacing $tu$, $tu'$, $tu''$, $t'u'$, $t'u''$, etc., with $(t,u)$, $(t,u')$,
$(t,u'')$, $(t',u')$, $(t',u'')$, etc.

Every point in space admits a trinomial
\[
(xt + x't' + x''t'')
\]
as its representation.

For all points lying in a given plane, an equation
\[
\lambda x + \lambda' x' + \lambda'' x'' = L
\]
then holds, where $\lambda$, $\lambda'$, $\lambda''$, $L$ denote fixed numbers. For a
plane through the three points $\mu t$, $\mu' t'$, $\mu'' t''$ we have
\[
\lambda\mu = \lambda'\mu' = \lambda''\mu'' = L.
\]
